% Zone „Das kennst du schon“ – aus bank/normalverteilung-und-sigma-regeln/zone.jsonl (zusammenbau v0.1)
\einheitenkopf[zone][Kennst du schon]{Das kennst du schon}
\setcounter{aufgabe}{0}

%% TODO Titel: Ich-kann-Satz fehlt in der Bank (Platzhalter: Kettenname) – Nr. 1
%% TODO Anweisung: ein Satz über der ersten Teilaufgabe fehlt in der Bank – Nr. 1
\begin{aufgabe}{Erwartungswert und Standardabweichung als Kenngrößen deuten}
\begin{teile}
\teil $X$ nimmt $1$ und $3$ je mit Wahrscheinlichkeit $0{,}5$ an. Erwartungswert? E(X) = \leerfeld
\teil $X$ nimmt $-2$, $0$ und $4$ mit den Wahrscheinlichkeiten $0{,}25$, $0{,}5$ und $0{,}25$ an. Erwartungswert? E(X) = \leerfeld
\end{teile}
\end{aufgabe}

%% TODO Titel: Ich-kann-Satz fehlt in der Bank (Platzhalter: Kettenname) – Nr. 2
%% TODO Anweisung: ein Satz über der ersten Teilaufgabe fehlt in der Bank – Nr. 2
\begin{aufgabe}{Flächeninhalt unter einem Graphen als Größe deuten}
\begin{teile}
\teil Die Zuflussrate ist $0{,}5$ Liter pro Minute von $t = 0$ bis $t = 2$ min. Fläche unter dem Graphen – was bedeutet sie? \leerfeld[l]
\teil Ein Zug fährt $1{,}5$ min lang mit $12$ m/s. Fläche unter dem Geschwindigkeitsgraphen? \leerfeld[m]
\end{teile}
\end{aufgabe}

%% TODO Titel: Ich-kann-Satz fehlt in der Bank (Platzhalter: Kettenname) – Nr. 3
%% TODO Anweisung: ein Satz über der ersten Teilaufgabe fehlt in der Bank – Nr. 3
\begin{aufgabe}{Binomialverteilung und ihr Säulendiagramm lesen}
\begin{teile}
\teil Lies ab: $P(X = 2)$?

\binomialverteilung{4}{0.5}
\leerfeld
\teil $n = 20$, $p = 0{,}3$: Erwartungswert und Standardabweichung?
\end{teile}
\end{aufgabe}

%% TODO Titel: Ich-kann-Satz fehlt in der Bank (Platzhalter: Kettenname) – Nr. 4
%% TODO Anweisung: ein Satz über der ersten Teilaufgabe fehlt in der Bank – Nr. 4
\begin{aufgabe}{Gegenereignis und Prozentrechnung}
\begin{teile}
\teil $P(A) = 0{,}3$: $P(\text{nicht } A)$? \leerfeld
\teil $P(X \le 3) = 0{,}812$: $P(X > 3)$? \leerfeld
\end{teile}
\end{aufgabe}

%% TODO Titel: Ich-kann-Satz fehlt in der Bank (Platzhalter: Kettenname) – Nr. 5
%% TODO Anweisung: ein Satz über der ersten Teilaufgabe fehlt in der Bank – Nr. 5
\begin{aufgabe}{Den Rechner für Verteilungswahrscheinlichkeiten einsetzen}
\begin{teile}
\teil $X$ binomialverteilt mit $n = 10$, $p = 0{,}5$: $P(X = 5)$ mit dem Rechner? \leerfeld
\teil $X$ binomialverteilt mit $n = 20$, $p = 0{,}3$: $P(X \le 5)$? \leerfeld
\end{teile}
\end{aufgabe}

%% TODO Titel: Ich-kann-Satz fehlt in der Bank (Platzhalter: Kettenname) – Nr. 6
%% TODO Anweisung: ein Satz über der ersten Teilaufgabe fehlt in der Bank – Nr. 6
\begin{aufgabe}{Flächeninhalt unter einem Graphen als Größe deuten – Fehler finden}
\begin{teile}
\teil Ich finde den Fehler: Die Zuflussrate ist $3$ Liter pro Minute, $6$ Minuten lang. Lea: „Es fließen $3$ Liter zu.“
\end{teile}
\end{aufgabe}

%% TODO Titel: Ich-kann-Satz fehlt in der Bank (Platzhalter: Kettenname) – Nr. 7
%% TODO Anweisung: ein Satz über der ersten Teilaufgabe fehlt in der Bank – Nr. 7
\begin{aufgabe}{Flächeninhalt unter einem Graphen als Größe deuten – selbst rechnen}
\begin{teile}
\teil Die Zuflussrate ist $2{,}5$ Liter pro Minute, $8$ Minuten lang. Wie viel Wasser fließt zu? \leerfeld[l]
\end{teile}
\end{aufgabe}
